\section{\sys}
\label{sec:method}

\begin{figure*}[t]
\centering
\footnotesize
\begin{tikzpicture}[
  >=Stealth,
  box/.style={draw=black!65, rounded corners=2pt, align=center, inner sep=2.5pt, fill=white, minimum height=0.58cm},
  skill/.style={box, minimum width=2.7cm, fill=blue!10},
  rulebox/.style={box, minimum width=2.7cm, fill=black!7},
  tool/.style={box, minimum width=2.7cm, fill=green!10},
  ctx/.style={box, minimum width=2.7cm, fill=yellow!12, text width=2.55cm},
  added/.style={ultra thick, draw=orange!85!black},
  loop/.style={box, text width=7.6cm, align=left, minimum height=0.82cm, fill=orange!7},
  head/.style={font=\footnotesize\bfseries}]
% harness before
\node[head] at (0,4.55) {Harness $H_{n-1}$ (one sentence)};
\node[ctx] (r0) at (0,3.80) {registry: 5 previous sentences + annotations};
\node[skill] (g0) at (0,2.95) {\skill{generate} \scriptsize(G)};
\node[skill] (c0) at (0,2.25) {\skill{critique} \scriptsize(C)};
\node[skill] (f0) at (0,1.55) {\skill{finalize} \scriptsize(G)};
\node[box, minimum width=2.7cm, fill=white, dashed] (t0) at (0,0.70) {tools: none (seed)};
\draw[->] (r0)--(g0); \draw[->] (g0)--(c0); \draw[->] (c0)--(f0);
\node[font=\scriptsize\itshape, text width=2.8cm, align=center] at (0,-0.15) {$H_0$: written by the builder from the guideline alone};
% loop
\node[head] at (6.0,4.55) {Round $n$ on training group $n$ (5 sentences, one document)};
\node[loop] (s1) at (6.0,3.80) {\textbf{1 Annotate} the group with $H_{n-1}$; keep every step's output.};
\node[loop] (s2) at (6.0,2.90) {\textbf{2 Judge} against gold: per-check units (umrjudge / timemljudge).};
\node[loop] (s3) at (6.0,2.00) {\textbf{3 Reflect} (builder): summarize problems, find their causes, propose operations on skills, steps and tools; may read resource files (PropBank, \ldots) and search the guideline.};
\node[loop] (s4) at (6.0,1.00) {\textbf{4 Admit} by machine checks: anchors, no test text, size budget, tool tests in a sandbox.};
\node[loop] (s5) at (6.0,0.10) {\textbf{5 Re-annotate} the group with $H_n$ (training score after) and \textbf{test} $H_n$ on held-out documents.};
\foreach \a/\b in {s1/s2,s2/s3,s3/s4,s4/s5} {\draw[->] (\a)--(\b);}
% harness after
\node[head] at (12.1,4.55) {Harness $H_n$};
\node[ctx] (r1) at (12.1,3.80) {registry: 5 previous sentences + annotations};
\node[skill] (g1) at (12.1,2.95) {\skill{generate} \scriptsize(G)\,$\pm$};
\node[skill] (c1) at (12.1,2.25) {\skill{critique} \scriptsize(C)\,$\pm$};
\node[skill,added] (n1) at (12.1,1.55) {new check step \scriptsize(C)};
\node[skill] (f1) at (12.1,0.85) {\skill{finalize} \scriptsize(G)\,$\pm$};
\node[tool,added] (t1) at (12.1,0.10) {tool: \texttt{propbank\_lookup}};
\draw[->] (r1)--(g1); \draw[->] (g1)--(c1); \draw[->] (c1)--(n1); \draw[->] (n1)--(f1);
% flows
\draw[->,thick] (g0.east) -- (s1.west);
\draw[->,thick] (s4.east) -- ++(0.5,0) |- (n1.west);
\draw[->,dashed] (t1.south) -- ++(0,-0.55) -| (0,-0.55) -- (0,-0.45);
\node[font=\scriptsize\itshape, fill=white, inner sep=1.5pt] at (6.0,-0.90) {$H_n$ is the starting point of round $n{+}1$; the builder never sees test documents or scores};
\end{tikzpicture}
\caption{\sys. Left: the seed harness annotates one sentence at a time; each step sees the previous five sentences with their annotations. Middle: one round of recursive self-improvement on a group of five gold sentences. Right: the next version; skills are edited ($\pm$), and steps and tools may be added (orange), e.g.\ a PropBank lookup the builder writes after it learns that frame files exist. G and C are the generator and critic models of a configuration.}
\label{fig:overview}
\end{figure*}


\subsection{Overview}
\label{sec:overview}

\sys turns an annotation guideline and a few annotated sentences into an annotator without training any model. The method is independent of the annotation scheme: a scheme supplies its guideline, an output format with a parser, a fine-grained judge, and optionally resource files; everything else is shared (Figure~\ref{fig:overview}).

\paragraph{Annotator.} An annotator is a harness $H=(S,T,G,R)$ run by language models: \emph{skills} $S$, the instructions of each model step; \emph{tools} $T$, typed Python functions a step may call; a \emph{chain} $G$, the order of the steps and the model that runs each; and \emph{rule steps} $R$, code that rewrites an annotation without a model. Each step receives the input text, the current annotation, and unresolved findings from earlier steps. A \emph{rewrite} step produces a complete annotation, a \emph{check} step identifies problems and proposes fixes for the next rewrite step, and a rule step runs code. The final annotation is the last one that parses. Every step receives the same output format; all other instructions are learned.

\paragraph{Seed.} The builder model reads the guideline and writes the seed harness $H_0$: three steps, \skill{generate} (writes the annotation), \skill{critique} (lists problems) and \skill{finalize} (rewrites the annotation with the findings), each a checklist of at most 8{,}000 characters. The builder sees no annotated data, and the seed has no tools.

\paragraph{Recursive self-improvement.} Training data come as groups of five consecutive sentences of one document. Round $n$ turns $H_{n-1}$ into $H_n$ with group $n$ (\S\ref{sec:rsi}): annotate the group, score it with the judge, let the builder summarize the problems, reflect on their causes and propose operations on skills, steps and tools, admit the operations that pass machine checks, and annotate the group again. Each $H_n$ is tested on held-out documents, which the builder never sees.

\paragraph{Judge.} The builder needs to know which part of the scheme each error belongs to. A judge aligns the predicted annotation with the gold one and scores a set of checks, each covering one part of the scheme (for example, node concepts, argument roles, attribute values, or one type of link). It returns \emph{units}, each with a verdict (correct, wrong, missing, extra), the sentence, and the gold and predicted elements. The judges we use are described in \S\ref{sec:benchmarks}.

\subsection{The harness and its improvement}
\label{sec:rsi}

\paragraph{Representation.} A harness version is a directory of plain files: \texttt{harness.yaml} (the chain: each step's name, kind, model and tools), one Markdown skill per model step, and one Python module per tool and per rule step. A version is immutable; round $n$ copies $H_{n-1}$ and edits the copy. A configuration names two models: the generator G runs \skill{generate} and \skill{finalize}, the critic C runs \skill{critique}, and a step the builder adds names which of the two runs it.

\paragraph{Resources.} The annotating steps see only their input. Tools read resource files through one module. The builder receives a manifest of these files (for example lexicons, lists of labels and the guideline itself) and may open them while it reflects. It receives no instruction about what to build: the seed has no tools, and every tool is the builder's own design.

\paragraph{One round.} (1) \emph{Annotate}: $H_{n-1}$ annotates the five sentences of group $n$; the sentences before the group are given with their gold annotations as context. The trace keeps every step's output and findings. (2) \emph{Judge}: the judge scores the group and lists the units that are not correct. (3) \emph{Reflect}: one builder call receives the output format, the harness (chain, all skills with their sizes, all tools), the resource manifest, the changes and training scores of earlier rounds (never a test score), and, for each sentence, its gold annotation, every step's output and its incorrect units. It returns \emph{problems} (error patterns with evidence and the step to fix), a \emph{reflection} (why they happen and what to change), \emph{operations}, and \emph{suggestions} that the operations cannot express. During the call the builder can list and read resource files and search the guideline. (4) \emph{Admit}: operations are applied in order to a copy of $H_{n-1}$ when they pass the machine checks. (5) \emph{Re-annotate} the group with $H_n$, and test $H_n$ on held-out documents.

\paragraph{Operations.} \texttt{edit\_skill} modifies, inserts or deletes text at an anchor that occurs exactly once; \texttt{rewrite\_skill} replaces a skill to consolidate it; \texttt{add\_step} inserts a check, rewrite or rule step after an existing step, and \texttt{delete\_step} removes one (at most five model steps); \texttt{add\_tool}, \texttt{replace\_tool} and \texttt{delete\_tool} manage tools, and \texttt{set\_step\_tools} assigns tools to steps.

\paragraph{Admission.} Machine checks decide admission; there is no reviewer model, so a second opinion cannot block a new tool or step. No text may copy a span of a test sentence (eight characters in scripts without word spaces, six words otherwise). A skill may hold at most 10{,}000 characters; an edit that would exceed this is rejected, and the builder is asked to consolidate with \texttt{rewrite\_skill}, which keeps the harness compact as rounds accumulate. A tool or rule may import only whitelisted modules (the resource module included), may use no file, evaluation or reflection primitives, must define one public function with type annotations and a docstring, and must pass at least three tests of its own in a separate process; a tool that fails returns once to the builder with the error. A new step follows an existing one, and the first step stays a rewrite step.
